\ExerciseSection

\begin{enumerate}
\def\labelenumi{\arabic{enumi})}
\tightlist
\item
  一个正实值函数 \(f\) （定义在 \(\mathbf{X} \times \mathbf{X}\) 上）是 \(\mathbf{X}\) 上的伪度量，当且仅当 \(f(x, x) = 0\) 对所有 \(x \in \mathbf{X}\) 成立，并且
\end{enumerate}

\[
f (x, y) \leqslant f (x, z) + f (y, z)
\]

对所有 \(x, y, z\) ∈ \(\mathbf{X}\) 成立 。

\begin{enumerate}
\def\labelenumi{\arabic{enumi})}
\setcounter{enumi}{1}
\tightlist
\item
  设 \(f\) 是 \(\mathbf{X} \times \mathbf{X}\) 到 \([0, +\infty]\) 中的一个映射。那么，由诸集 \(\overline{f}^{1}([0, a])\) （其中 \(a\) 取遍所有实数 \(>0\) ）组成的族，构成 \(\mathbf{X}\) 上某个一致结构的围域基本系，当且仅当 \(f\) 满足下列条件：\(a) f(x, x) = 0\) 对所有 \(x \in \mathbf{X}\) 成立；\(b)\) 对每个 \(a > 0\) ，存在 \(b > 0\) ，使得关系 \(f(x, y) \leqslant b\) 蕴含 \(f(y, x) \leqslant a\) ；\(c)\) 对每个 \(a > 0\) ，存在 \(c > 0\) ，使得关系 \(f(x, z) \leqslant c\) 与 \(f(z, y) \leqslant c\) 蕴含 \(f(x, y) \leqslant a\) 。
\end{enumerate}

条件 \(b)\) 与 \(c)\) 特别在下述情形得到满足：存在 \(\varphi\) ——一个 \(I = [0, +\infty]\) 到其自身上的连续映射，它在点 \(0\) 处取值零——以及 \(\psi\) ——一个 \(I \times I\) 到 \(I\) 中的连续映射，它在点 \((0, 0)\) 处取值零——使得恒有

\[
f (y, x) \leqslant \varphi (f (x, y)) \quad \text { and } \quad f (x, y) \leqslant \psi (f (x, z), f (z, y)).
\]

\begin{enumerate}
\def\labelenumi{\arabic{enumi})}
\setcounter{enumi}{2}
\tightlist
\item
  设 X 是一个拓扑空间，C 是 X 的拓扑，\((f_{t})\) 是 X 上伪度量的一个饱和族（第 2 号），U 是由族 \((f_{t})\) 定义的一致结构。
\end{enumerate}

\begin{enumerate}
\def\labelenumi{\alph{enumi})}
\item
  由 \(\mathcal{U}\) 诱导的拓扑粗于 \(\mathfrak{C}\) ，当且仅当诸 \(f_{t}\) 在 \(\mathbf{X} \times \mathbf{X}\) 上连续（关于拓扑 \(\mathfrak{C}\) 与其自身的乘积）。
\item
  由 \(\mathcal{U}\) 诱导的拓扑细于 \(\mathcal{C}\) ，当且仅当对每个 \(x_0 \in \mathbf{X}\) 与每个邻域 \(\mathbf{V}\) ——它包围 \(x_0\) （按拓扑 \(\mathcal{C}\) ）——，存在一个指标 \(\iota\) 与一个实数 \(a > 0\) ，使得 \(f_{\iota}(x_0, x) \geqslant a\) 对所有 \(x \in \mathbb{C}\mathbf{V}\) 成立 。
\end{enumerate}

\begin{enumerate}
\def\labelenumi{\arabic{enumi})}
\setcounter{enumi}{3}
\tightlist
\item
  设 X 是一个非豪斯多夫的可一致化空间，R 是关系`` \(y \in \overline{\{x\}}\) ''（在 X 的两点 \(x, y\) 之间）。
\end{enumerate}

\begin{enumerate}
\def\labelenumi{\alph{enumi})}
\item
  证明 R 是 X 上的一个等价关系，X 到豪斯多夫空间中的每个连续映射都与关系 R 相容（《集合论》，R，§ 5，第 7 号），并且商空间 X/R 是完全正则的。
\item
  若 \(\mathcal{U}\) 是任一与 X 的拓扑相容的一致结构，则与 \(\mathcal{U}\) 相伴的豪斯多夫一致结构（第 II 章，§ 3，第 9 号）定义在商空间 X/R 上，且与 X/R 的拓扑相容。拓扑空间 X/R 称为与可一致化空间 X 相伴的完全正则空间。
\end{enumerate}

\begin{enumerate}
\def\labelenumi{\arabic{enumi})}
\setcounter{enumi}{4}
\tightlist
\item
  设 X 是一个可一致化空间。证明 X 上所有在 \(X \times X\) 上连续的伪度量构成的族定义了 X 上一个与 X 的拓扑相容的一致结构。这个一致结构 \(U_{0}\) 称为 X 上的万有一致结构；它是与 X 的拓扑相容的所有一致结构中最细的一个。若 Y 是任一一致空间，且 \(f: X \to Y\) 是连续映射，则 f 关于 X 上的万有一致结构是一致连续的；对于 X 上与其拓扑相容的任何其他一致结构，情形并非如此。若存在 X 上的一致结构 U ，它与 X 的拓扑相容，且使 X 赋有该一致结构后为完备空间，则 X 关于任何一致结构 \(U'\) （它比 U 更细且比 \(U_{0}\) 更粗）也是完备空间。（注意，柯西滤子关于 \(U'\) 与关于 \(U_{0}\) 是同样的。）
\end{enumerate}

¶ 6) a) 设 X 是一个完全正则空间，U 是一个与 X 的拓扑相容的一致结构。设 U* 是 X 上使所有关于 U 一致连续的 X 到 {[}0, 1{]} 中的映射都一致连续的最粗一致结构。证明一致结构 U* 是豪斯多夫的且与 X 的拓扑相容，并且 X 关于 U* 是预紧的。

\begin{enumerate}
\def\labelenumi{\alph{enumi})}
\setcounter{enumi}{1}
\tightlist
\item
  设 K 是一个紧空间。证明每个关于 U 一致连续的映射 \(f: X \to K\) 关于 \(U^{*}\) 也一致连续（注意，K 上唯一的一致结构是使 K 到区间 {[}0, 1{]} 中的所有连续映射都一致连续的最粗一致结构）。由此证明 \(U^{*}\) 是 X 上比 U 更粗且使 X 关于它为预紧的一致结构中最细的一个。
\end{enumerate}

\begin{enumerate}
\def\labelenumi{\arabic{enumi})}
\setcounter{enumi}{6}
\tightlist
\item
  设 X 是一个完全正则空间。若 \(U_{0}\) 表示 X 上的万有一致结构（习题 5），则 \(U_{0}^{*}\) （习题 6）是 X 上使所有 X 到 {[}0, 1{]} 中的连续映射都一致连续的最粗一致结构。设 \(\beta X\) 表示关于一致结构 \(U_{0}^{*}\) 完备化 X 所得的紧空间。\(\beta X\) 称为 X 的斯通–切赫紧化。
\end{enumerate}

\begin{enumerate}
\def\labelenumi{\alph{enumi})}
\item
  证明 X 到紧空间 K 中的每个连续映射都可以扩张为 \(\beta X\) 到 K 中的连续映射（参看《集合论》，第 IV 章，§ 3）。
\item
  设 \(f\) 是 \(\mathbf{X}\) 到一个稠密子集 \(\mathbf{X}'\) （紧空间 \(\mathbf{K}\) 的）上的同胚，\(\overline{f}\) 是连续映射 \(\beta \mathbf{X} \to \mathbf{K}\) （它扩张 \(f\) ）。证明 \(\overline{f}(\beta \mathbf{X} - \mathbf{X}) = \mathbf{K} - \mathbf{X}'\) 。{[}假设存在一点 \(a \in \beta \mathbf{X} - \mathbf{X}\) 使 \(\overline{f}(a) = f(b)\) ，其中 \(b \in \mathbf{X}\) ；考察连续映射 \(g: \beta \mathbf{X} \to [0, 1]\) ，它满足 \(g(a) = 1\) 与 \(g(b) = 0\) ，并令 \(P\) 表示所有 \(x \in \beta X\) （满足 \(g(x) \geqslant 1/2\) ）所成的集。存在连续映射 \(h: K \to [0, 1]\) ，使 \(h(f(b)) = 0\) 且 \(h(f(x)) = 1\) 对所有 \(x \in P\) 成立；证明关系 \(h(f(a)) = 0\) 与 \(a\) 属于 \(P\) 在 \(\beta X\) 中的闭包这一事实相矛盾。{]}
\end{enumerate}

\begin{enumerate}
\def\labelenumi{\arabic{enumi})}
\setcounter{enumi}{7}
\tightlist
\item
  设 X 是一个完全正则空间，\(\beta X\) 是它的斯通–切赫紧化。
\end{enumerate}

\begin{enumerate}
\def\labelenumi{\alph{enumi})}
\item
  X 上的滤子 \(\mathfrak{F}\) 称为完全正则滤子，如果 \(\mathfrak{F}\) 有一个由开集组成的基 \(\mathfrak{B}\) ，使得对每个集 \(A \in \mathfrak{B}\) ，存在含于 A 的集 \(B \in \mathfrak{B}\) ，以及一个连续映射 \(f: X \to [0, 1]\) ，它在 B 上等于 0 、在 \(\mathbb{C}A\) 上等于 1 。完全正则滤子称为极大的，如果不存在严格更细的完全正则滤子。证明：若 \(\mathfrak{F}\) 是 E 上任一完全正则滤子，则存在一个比 \(\mathfrak{F}\) 更细的极大完全正则滤子（用佐恩引理）。
\item
  完全正则滤子 \(\mathfrak{F}\) 是极大的，当且仅当：对 X 的每对开集 A, B ，其中 \(B \subset A\) 且存在 X 到 \([0, 1]\) 中的连续映射在 B 上等于 0 、在 \(\complement A\) 上等于 1 ，则或者 \(A \in \mathfrak{F}\) ，或者 \(A \notin \mathfrak{F}\) 而 \(\mathfrak{F}\) 中有一个不与 B 相交的集（若 \(\mathfrak{F}\) 的所有集都与 B 相交，考察由 \(\mathfrak{F}\) 的诸集与诸集 \(\overline{f}([0, a[)\) （其中 \(a \in ]0, 1[\) ）生成的滤子）。
\item
  证明每个极大完全正则滤子 \(\mathfrak{F}\) 都是 \(\mathbf{X}\) 上关于由 \(\beta\mathbf{X}\) 的一致结构诱导的一致结构的柯西滤子{[}证明每个连续映射 \(f: \mathbf{X} \to [\mathrm{0}, \mathrm{1}]\) 关于 \(\mathfrak{F}\) 有极限（用反证法，并利用 \(b)\) ）{]}。
\end{enumerate}

\(d)\) 两个不同的极大完全正则滤子 \(\mathfrak{F},\mathfrak{F}^{\prime}\) 不能收敛到 \(\beta\mathbf{X}\) 的同一点{[}注意，由 \(b)\) ，存在开集 \(\mathbf{A}\in\mathfrak{F}\) 与开集 \(\mathbf{A}^{\prime}\in\mathfrak{F}^{\prime}\) ，使 \(\mathbf{A}\cap\mathbf{A}^{\prime}=\varnothing\) ；然后利用公理 \((\mathrm{O}_{\mathbf{IV}})\) 进行反证{]}。

\begin{enumerate}
\def\labelenumi{\alph{enumi})}
\setcounter{enumi}{4}
\tightlist
\item
  任一点 \(x_{0} \in \beta X\) 的邻域滤子在 X 上的迹是一个极大完全正则滤子 B （用反证法，并利用完全正则滤子的定义以及 \(\beta X\) 中一点的邻域滤子的定义）。
\end{enumerate}

\begin{enumerate}
\def\labelenumi{\arabic{enumi})}
\setcounter{enumi}{8}
\tightlist
\item
  \begin{enumerate}
  \def\labelenumii{\alph{enumii})}
  \tightlist
  \item
    设 X 是一个拓扑空间，\(\mathcal{G}\) 是它的拓扑，\((f_{\iota})_{\iota \in \mathbb{I}}\) 是 X 到 K = {[}0, 1{]} 中的一族映射，它们在拓扑 \(\mathcal{G}\) 中连续。设 \(\mathcal{G}_0\) 是 X 上使所有 \(f_{\iota}\) 都连续的最粗拓扑。证明 \(\mathcal{G}_0 = \mathcal{G}\) ，只要诸集 \(\overline{f}_{\iota}^{1}([0, a[)\) （其中 \(\iota \in I\) 且 \(a \in ]0, 1[,\) ）组成的族是 \(\mathcal{G}\) 的一个子基（第 I 章，§ 2，第 3 号）。此时空间 X 是可一致化的，并且若它是豪斯多夫的，它就同胚于方体 K \(^{\mathbf{I}}\) 的一个子空间。
  \end{enumerate}
\end{enumerate}

\begin{enumerate}
\def\labelenumi{\alph{enumi})}
\setcounter{enumi}{1}
\tightlist
\item
  假设族 \((f_{t})\) 是 X 到 K 中的所有连续映射（关于 C）构成的族。证明：若 Y 是任一紧空间，则每个关于 C 连续的 X 到 Y 中的映射关于 \(C_{0}\) 也连续（把 Y 嵌入一个方体）。拓扑 C 是可一致化的，当且仅当 \(C_{0} = C\) 。
\end{enumerate}

\begin{enumerate}
\def\labelenumi{\arabic{enumi})}
\setcounter{enumi}{9}
\tightlist
\item
  设 X 是一个完全正则空间，K 是 X 的一个紧子集，V 是 K 在 X 中的一个邻域。
\end{enumerate}

\begin{enumerate}
\def\labelenumi{\alph{enumi})}
\item
  证明存在 X 到 {[}0, 1{]} 中的连续映射，它在 K 上等于 1 ，在 \(\complement V\) 上等于 0 {[}利用 \((\mathrm{O}_{\mathbf{IV}})\) ，用有限多个适当选取的邻域覆盖 K{]}。
\item
  设 \(X'\) 是把 K 的所有点等同起来所得到的 X 的商空间。证明 \(X'\) 是完全正则的。
\end{enumerate}

¶ 11) 设 X 是一个完全正则空间。X 中两个闭集 A, B 称为完全分离的，如果存在 X 到 {[}0, 1{]} 中的连续映射 f ，使得在 A 上 \(f(x) = 0\) ，在 B 上 \(f(x) = 1\) 。a) 设 \(\beta X\) 是 X 的斯通–切赫紧化。证明 X 中两个闭集 A 与 B 完全分离，当且仅当它们在 \(\beta X\) 中的闭包不相交{[}利用习题 7 a) 与 10 a){]}。

\begin{enumerate}
\def\labelenumi{\alph{enumi})}
\setcounter{enumi}{1}
\tightlist
\item
  设 \(h\) 是 \(\mathbf{X}\) 到一个稠密子集 \(\mathbf{X}'\) （紧空间 \(\mathbf{K}\) 的）上的同胚，\(\overline{h}\) 是连续映射 \(\beta \mathbf{X} \to \mathbf{K}\) （它扩张 \(h\) ）。假设对每对完全分离的闭子集 \(\mathbf{A}, \mathbf{B}\) （属于 \(\mathbf{X}\) ），\(\overline{h}(\mathbf{A})\) 与 \(h(\mathbf{B})\) 在 \(\mathbf{K}\) 中的闭包不相交。证明 \(\overline{h}\) 是 \(\beta \mathbf{X}\) 到 \(\mathbf{K}\) 上的一个同胚。（证明 \(\overline{h}\) 是单射：若 \(a, b\) 是 \(\beta \mathbf{X}\) 中满足 \(\overline{h}(a) = \overline{h}(b)\) 的两个不同点，考察连续映射 \(f: \beta \mathbf{X} \to [\mathbf{0}, \mathbf{1}]\) ，它满足 \(f(a) = 0\) 、\(f(b) = 1\) ，并考察集 \(X \cap \overline{f}[0, 1/3]\) 与 \(X \cap \overline{f}[2/3, 1]\) 。）
\end{enumerate}

¶ 12) 设 X 是一个完全正则空间，βX 是它的斯通–切赫紧化。

\begin{enumerate}
\def\labelenumi{\alph{enumi})}
\item
  设 \(f\) 是 \(\beta X\) 上的一个有限连续实值函数，使 \(f(x) > 0\) 对所有 \(x \in X\) 成立，且使集 \(A = \overline{f}^1(0) \subset \beta X - X\) 非空。那么存在点的序列 \((a_n)\) （属于 \(X\) ），使数列 \(\lambda_n = f(a_n) > 0\) 严格递减，且 \(\lim_{n \to \infty} \lambda_n = 0\) 。对每个整数 \(n > 0\) ，设 \(I_n\) 是以 \(\lambda_n\) 为中心的开区间（在 \(R_+^*\) 中），使得闭区间 \(\overline{I}_n\) 两两不相交，并令 \(M_n = \overline{f}^1(I_n) \cap X\) 。对每个子集 \(H\) ⊂ \(N\) ，令 \(M_H\) 为诸 \(M_n\) （其中 \(n \in H\) ）的并；对每个滤子 \(\mathfrak{F}\) （在 \(N\) 上且比弗雷歇滤子更细），令 \(\mathfrak{F}'\) 表示滤子（在 \(X\) 上），它以诸 \(M_H(H \in \mathfrak{F})\) 为基。证明 \(\mathfrak{F}'\) 是完全正则滤子（习题 8），并且若 \(\mathfrak{F}_1, \mathfrak{F}_2\) 是两个不同的超滤子（在 \(N\) 上），其中每个都比弗雷歇滤子更细，则 \(X\) 上不存在比 \(\mathfrak{F}_1', \mathfrak{F}_2'\) 中每个都更细的滤子。由此推出 Card \((A) \geqslant 2^{2\text{Card}(N)}\) {[}利用习题 8 d) 与第 I 章，§ 4，习题 5{]}。
\item
  假设 X 是无限且离散的。证明由 \(\beta\mathbf{X}\) 的一致结构在 X 上诱导的一致结构是有限分划的一致结构（第 II 章，§ 2，第 2 号及 § 4，习题 12）{[}利用习题 11 b){]}。由此证明 Card \((\beta\mathbf{X}) = 2^{2^{\operatorname{Card}(\mathbb{N})}}\) （参看第 I 章，§ 4，习题 5）。证明 X 上存在连续实值函数 \(f \geqslant 0\) ，使 \(\overline{f}^{1}(0)\) 非空且含于 \(\beta\mathbf{X} - \mathbf{X}\) 。
\item
  在与 b) 相同的假设下，证明：若 A 是 \(\beta X\) 的一个无限闭子集，则 \(\operatorname{Card}(A) \geqslant 2^{2\operatorname{Card}(N)}\) 。{[}首先证明存在点的无限序列 \((a_{n})\) （属于 A ），且对每个 n 存在邻域 \(V_{n}\) （\(a_{n}\) 在 \(\beta X\) 中的），诸 \(V_{n}\) 两两不相交；由此推出定义在诸 \(a_{n}\) 所成的集 D 上的每个有界实值函数 f 都可以扩张为 \(\beta X\) 上的连续函数；为此，考察定义在 X 上的实值函数 g ，使 \(f(a_{n})\) 是它在 \(X \cap V_{n}\) 每点处的值。于是断言 \(\overline{D} \subset A\) 同胚于 D 的斯通–切赫紧化（参看 § 4，习题 17）。{]}
\end{enumerate}

\begin{enumerate}
\def\labelenumi{\arabic{enumi})}
\setcounter{enumi}{12}
\item
  设 G 是一个局部紧但非紧的拓扑群。证明：在乘积空间 \(H = G \times G\) 上，由其斯通–切赫紧化 \(\beta H\) 的一致结构诱导的一致结构，严格细于由 \((\beta G) \times (\beta G)\) 的一致结构诱导的一致结构。{[}假设结论不真，把习题 7 a) 应用于 H 到 G 中的映射 \((x, y) \to xy^{-1}\) ，证明 G 将同构于一个紧群的子群；然后应用第 III 章，§ 3，第 3 号，命题 4 的推论 1 。{]}
\item
  若拓扑空间 X 的每一点都有一个闭邻域，它是 X 的可一致化子空间，证明 X 是可一致化的。
\end{enumerate}

¶ 15) a) 设 X 是一个局部紧空间，\(\Phi\) 是所有满足下列条件的连续映射 \(g: \mathbf{X} \to [\mathfrak{0}, \mathfrak{1}]\) 组成的集：g 在某个紧集（依赖于 \(g\) ）的余集上为零。设 \(\mathcal{U}_1\) 是 X 上使所有映射 \(g \in \Phi\) 都一致连续的最粗一致结构。证明 \(\mathcal{U}_1\) 与 X 的拓扑相容，X 的完备化 \(\hat{\mathbf{X}}\) （关于 \(\mathcal{U}_1\) ）是紧的，且 \(\mathbf{X} - \hat{\mathbf{X}}\) 至多由一点组成（证明：若 X 不是紧的，则 X 的相对紧子集的余集所成的滤子是关于 \(\mathcal{U}_1\) 的柯西滤子）。

\begin{enumerate}
\def\labelenumi{\alph{enumi})}
\setcounter{enumi}{1}
\item
  证明 \(\mathfrak{U}_1\) 是与 \(\mathbf{X}\) 的拓扑相容的最粗一致结构。
\item
  反之，设 X 是一个完全正则空间，使得与 X 的拓扑相容的一致结构组成的集有一个最粗元 \(U_{1}\) 。证明 X 是局部紧的。（利用习题 6 证明 X 关于 \(U_{1}\) 是预紧的；然后证明 X 在其关于 \(U_{1}\) 的完备化中的余集不可能多于一点。）
\item
  假设 X 是局部紧且 \(\sigma\) 紧的，因而是递增序列 \((\mathbf{U}_n)\) （由相对紧开集组成，且满足 \(\overline{\mathbf{U}}_n \subset \mathbf{U}_{n+1}\) ）之并（第 I 章，§ 9，第 9 号，命题 15）。证明存在一个连续实值函数 \(f\) （定义在 \(\mathbf{X}\) 上），使得 \(f(x) \leqslant n\) 对 \(x \in \overline{\mathbf{U}}_n\) 成立，且 \(f(x) \geqslant n\) 对 \(x \in \mathbb{C} \overline{\mathbf{U}}_n\) 成立{[}参看习题 10 a){]}。设 \(\mathcal{U}_2\) 是 \(\mathbf{X}\) 上使 \(f\) 与所有函数 \(g \in \Phi\) 都一致连续的最粗一致结构。证明 \(\mathcal{U}_2\) 与 \(\mathbf{X}\) 的拓扑相容，且存在一个围域 \(\mathbf{V}\) （属于 \(\mathcal{U}_2\) ），使得 \(\mathbf{V}(x)\) 在 \(\mathbf{X}\) 中相对紧对所有 \(x \in \mathbf{X}\) 成立（参看第 II 章，§ 4，习题 9）。
\end{enumerate}

¶ 16) 设 X 是一个完备的豪斯多夫一致空间，U 是 X 的一致结构。

\begin{enumerate}
\def\labelenumi{\alph{enumi})}
\tightlist
\item
  设 A 是 X 的一个子集，它是一个闭集序列 \((\mathbf{F}_{n})\) 之并，且 \(x \notin A\) 。证明 X 上存在连续的有限实值函数 \(f \geqslant 0\) ，使 \(f(x) = 0\) 且 \(f(y) > 0\) 对所有 \(y \in A\) 成立。（若 \(g_{n}\) 是 X 到 {[}0, 1{]} 中的连续映射，满足 \(g_{n}(x) = 0\) 且 \(g_{n}(y) = 1\) 对所有 \(y \in F_{n}\) 成立，考察函数
\end{enumerate}

\[
f (x) = \sum_ {n = 0} ^ {\infty} g _ {n} (x) / 2 ^ {n}.
\]

\begin{enumerate}
\def\labelenumi{\alph{enumi})}
\setcounter{enumi}{1}
\tightlist
\item
  设 \(h \geqslant 0\) 是 \(\mathbf{X}\) 上的连续实值函数，\(\mathbf{V}\) 是开集 \(\overline{h}^{1}(]0, +\infty[)\) 。设 \(\mathcal{U}'\) 是 \(\mathbf{V}\) 上由 U 诱导的一致结构与由伪度量
\end{enumerate}

\[
r (x, y) = \left| \frac {\mathrm{1}}{h (x)} - \frac {\mathrm{1}}{h (y)} \right|.
\]

定义的一致结构的最小上界。证明 V 关于 \(U'\) 是完备的。

\begin{enumerate}
\def\labelenumi{\alph{enumi})}
\setcounter{enumi}{2}
\tightlist
\item
  设 B 是 X 的一个子集，它是一族集 \((\mathbf{A}_{\lambda})\) 之交，其中每个集都是可数闭集族之并。证明 B 上存在一个一致结构，它与 X 的拓扑在 B 上诱导的拓扑相容，且 B 关于这个一致结构是完备的{[}利用 b){]}。
\end{enumerate}

\begin{enumerate}
\def\labelenumi{\arabic{enumi})}
\setcounter{enumi}{16}
\item
  设 X 是一个拓扑空间，它的每个点 \(x \in X\) 都有一个既开又闭的邻域组成的基本邻域系。证明 X 是可一致化的。
\item
  设 X 是一个可数的完全正则空间。证明对每个 \(x \in X\) ，x 的既开又闭的邻域构成 x 的一个基本邻域系（参看 § 6，第 4 号）。
\item
  设 X 是一个局部紧空间。证明 X 上每个下半连续的函数 \(f \geqslant 0\) 都是一族连续函数 \(\geqslant 0\) 的上包络，这族函数中的每一个都在某个紧集的余集上为零。
\item
  \begin{enumerate}
  \def\labelenumii{\alph{enumii})}
  \tightlist
  \item
    设 X 是一个完全正则空间，a 是 X 的一点。那么，存在 X 上的连续函数 \(f \geqslant 0\) ，使 \(f(a) = 0\) 且 \(f(x) > 0\) 对每个 \(x \neq a\) 成立，当且仅当存在 a 在 X 中的邻域序列 \((\mathbf{V}_n)\) ，使 \(\bigcap_{n} V_n = \{a\}\) {[}参看习题 16 a){]}。
  \end{enumerate}
\end{enumerate}

\begin{enumerate}
\def\labelenumi{\alph{enumi})}
\setcounter{enumi}{1}
\tightlist
\item
  设 X 是一个完全正则空间，它的每个点都有可数的基本邻域系。证明在斯通–切赫紧化 \(\beta\mathrm{X}\) 中，集 X 等于 \(\beta\mathrm{X}\) 中具有可数基本邻域系的点所成的集{[}利用 a) 与习题 12{]}。设 \(\mathbf{X}'\) 是另一个完全正则空间，它的每个点都有可数的基本邻域系。证明：若斯通–切赫紧化 \(\beta\mathrm{X}\) 与 \(\beta\mathrm{X}'\) 同胚，则 X 与 \(\mathbf{X}'\) 同胚。
\end{enumerate}

¶ 21) a) 设 X 是拓扑空间。证明下列两个性质等价：α) X 上每个连续有限实值函数都有界；β) X 上每个有界的连续实值函数都达到它的界。若 X 具有这两个性质，则称 X 为伪紧的。若 X 伪紧，且 f 是 X 到拓扑空间 \(X'\) 中的任一连续映射，则 \(f(X)\) 伪紧。

\begin{enumerate}
\def\labelenumi{\alph{enumi})}
\setcounter{enumi}{1}
\item
  考虑拓扑空间 X 的下列两个性质：\(\gamma\) 若 \((\mathbf{U}_n)\) 是 X 的任一可数开覆盖，则 X 是有限个闭包 \(\overline{\mathbf{U}}_n\) 的并；\(\delta\) X 上每个由开集组成的可数滤子基至少有一个触点。证明 \(\gamma\) 与 \(\delta\) 等价，且它们蕴含 X 伪紧。特别地，每个绝对闭的拓扑空间（第 I 章，§ 9，习题 19）都是伪紧的。若 X 具有性质 \(\gamma\) 与 \(\delta\) ，则 X 的每个作为 X 的开集之闭包的子空间也具有这两性质{[}参看 § 4，习题 26 b){]}。
\item
  考虑拓扑空间 \(\mathbf{X}\) 的下列性质：\(\zeta)\) \(\mathbf{X}\) 的每个局部有限开覆盖都是有限的。证明 \((\zeta) \Longrightarrow \gamma)\) 。{[}若 \((\mathbf{U}_n)\) 是 \(\mathbf{X}\) 的开子集组成的一个递增序列，诸子集覆盖 \(\mathbf{X}\) 且满足 \(\mathbf{U}_{n+1} \notin \overline{\mathbf{U}}_n\) ，考虑由诸集 \(\mathbf{U}_{n+1} \cap \mathbb{C}\overline{\mathbf{U}}_n\) 以及某个序列 \((a_n)\) 的余集所组成的覆盖，该序列满足
\end{enumerate}

\[
a _ {n + 1} \in \mathrm{U} _ {n + 1} \cap \left[ \overline {{\mathrm{U}}} _ {n}. \right]
\]

若 X 是正则的，证明 \(\gamma\Longrightarrow\zeta\) 。{[}设 \((\mathrm{U}_{\alpha})\) 是 X 的一个无限的局部有限开覆盖。用归纳法定义指标的序列 \((\alpha_{n})\) 、X 的点的序列 \((x_{n})\) ，并对每个 \(x_{n}\) 定义两个开邻域 \(V_{n}\) 与 \(W_{n}\) （\(x_{n}\) 的），使得：(i) \(\overline{V}_{n}\subset W_{n}\) ，\(W_{n}\subset U_{\alpha_{n}}\) ，且 \(W_{n}\) 只与有限个集 \(U_{\alpha}\) 相交；(ii) \(U_{\alpha_{n}}\) 不与任何 \(W_{k}\) （指标 k\textless n 者）相交。现在考虑由诸 \(W_{n}\) 与诸 \(\overline{V}_{n}\) 的并的余集所组成的覆盖。{]}

\begin{enumerate}
\def\labelenumi{\alph{enumi})}
\setcounter{enumi}{3}
\item
  在完全正则空间 X 中，性质 \(\alpha\) 、\(\beta\) 、\(\gamma\) 、\(\delta\) ) 与 \(\zeta\) 都等价，并且都与下列性质等价：\(\theta\) X 在任何与它的拓扑相容的一致结构中都是预紧的。{[}为证 \(\alpha) \Longrightarrow \zeta\) )，注意：若 \((\mathbf{U}_n)\) 是 X 的一个可数无限的局部有限开覆盖，且 \(a_n \in \mathbf{U}_n\) ，则 X 上存在连续实值函数 \(f\) ，使得 \(f(a_n) = n\) 对每个 \(n\) 成立{[}利用 \((\mathrm{O}_{\mathrm{IV}})]\) 。为证 \(\theta) \Longrightarrow \alpha\) )，利用习题 5，并注意在预紧空间中每个一致连续实值函数都有界。为证 \(\gamma) \Longrightarrow \theta\) )，注意：若 X 关于某个一致结构 \(\mathcal{U}\) 不是预紧的，则存在 \(\mathcal{U}\) 的一个对称围域 V 与 X 的点的序列 \((x_n)\) ，使得诸集 \(\mathrm{V}(x_n)\) 中任何两个都不相交。{]}
\item
  若 X 是完全正则的伪紧空间，且关于某个与它的拓扑相容的一致结构是完备的，则 X 是紧的。
\item
  在完全正则的伪紧空间 X 上，万有一致结构（习题 5）由 X 的斯通–切赫紧化的一致结构所诱导，并且是使 X 到 {[}0, 1{]} 的所有连续映射都一致连续的、唯一与 X 的拓扑相容的一致结构。
\end{enumerate}

\begin{enumerate}
\def\labelenumi{\arabic{enumi})}
\setcounter{enumi}{21}
\tightlist
\item
  设 X 是豪斯多夫一致空间，Y 是 X 的闭子集，f 是 Y 上有界的一致连续实值函数。证明 f 有一个扩张 \(\overline{f}\) 到 X 上，它在 X 上有界且一致连续。{[}我们可以假设 \(f(Y) \subset [0, 1]\) 。对每个二进数 \(r = k/2^{n} \in [0, 1]\) ，设 \(A(r)\) 是由所有 \(x \in Y\) （满足 \(f(x) \leqslant r\) ）组成的集，并设 \(B(r) = A(r) \cup (X - Y)\) 。用归纳法 * （按 § 4 第 1 号定理 1 的证明方式） * 定义一个开集 \(U(r)\) （对每个二进数 \(r \in [0, 1]\) ），使得：只要 r 与 \(r'\) 是两个二进数且 \(r < r'\) ，就存在 X 的一致结构的一个围域 V，使得
\end{enumerate}

\[
\mathbf {V} (\mathbf {A} (r)) \subset \mathbf {U} (r ^ {\prime}), \quad \mathbf {V} (\mathbf {U} (r)) \subset \mathbf {U} (r ^ {\prime}), \quad \mathbf {V} (\mathbf {U} (r)) \subset \mathbf {B} (r ^ {\prime}).
\]

然后定义 \(\overline{f}(x)\) 为诸二进数 \(r\) ——满足 \(x \in \mathrm{U}(r).]\) 者——的下确界。

